\documentclass[11pt]{article}
\usepackage[utf8]{inputenc}
\usepackage{amsmath, amssymb, amsthm}
\usepackage[margin=1in]{geometry}

\theoremstyle{definition}
\newtheorem{definition}{Definition}
\newtheorem{theorem}{Theorem}

\theoremstyle{remark}
\newtheorem{remark}{Remark}

\DeclareMathOperator{\tr}{tr}

\title{Lecture 4: Stochastic Integration (continued)}
\date{06.11.2025}
\author{Tien Anh Vu}
\begin{document}
\maketitle

%% =============================
%% §4 Stochastic integration (continued)
%% =============================
\section*{\S 4 Stochastic Integration (continued)}

%% =============================
%% Wiener--Itô isometry (recap)
%% =============================
\subsection*{Wiener--It\^o isometry (recap)}

We use the following notation for the stochastic integral. For an elementary (step) integrand
$\Phi \in \mathcal{E}$, define the (It\^o) stochastic integral process
\[
    I(\Phi)_t
    \;:=\;
    \int_0^t \Phi(s)\,dW(s),
    \qquad t\in[0,T],
\]
which is an $H$-valued continuous square-integrable martingale. Thus $I$ is the map
\[
    I \colon \mathcal{E} \longrightarrow \mathcal{M}_T^2,
    \qquad
    \Phi \longmapsto \bigl(I(\Phi)_t\bigr)_{t\in[0,T]}.
\]

Recall the Wiener--It\^o isometry from Step~2:
\[
    \bigl\|I(\Phi)\bigr\|_{\mathcal{M}_T^2}^2
    \;=\;
    \|\Phi\|_{\mathcal{E}}^2,
    \qquad
    I(\Phi) \in \mathcal{M}_T^2,\quad
    \Phi \in \mathcal{E}.
\]

\begin{remark}
Here $\mathcal{E}$ denotes the space of elementary (step) processes from Step~1, equipped with the norm
\[
    \|\Phi\|_{\mathcal{E}}^2 := E\!\bigl(\int_0^T \|\Phi_s \circ \sqrt{Q}\|_{L_2(U,H)}^2\,ds\bigr),
\]
and $\mathcal{M}_T^2$ is the space of continuous square-integrable $H$-valued martingales on $[0,T]$.
The isometry says that the stochastic integral $I \colon \mathcal{E} \to \mathcal{M}_T^2$ is a \textbf{norm-preserving} (isometric) linear map: the $L^2$-norm of the integral equals the norm of the integrand.
This is the starting point for extending $I$ beyond elementary processes by a density/completion argument.
\end{remark}

%% =============================
%% Step 3: Closure / extension of I
%% =============================
\subsection*{Step 3: Extension to the closure $\overline{\mathcal{E}}$}

\paragraph{What is the closure $\overline{\mathcal{E}}$\,?}

We want to extend $I$ from the elementary processes $\mathcal{E}$ to a larger space $\overline{\mathcal{E}}$.

\begin{remark}
Since $I \colon \mathcal{E} \to \mathcal{M}_T^2$ is an isometry, it extends uniquely by continuity to the \textbf{closure} of $\mathcal{E}$ in the $\|\cdot\|_{\mathcal{E}}$-norm.
The question is: what does this closure look like concretely?
The answer involves the \textbf{predictable $\sigma$-algebra}, which identifies the ``right'' class of integrands.
\end{remark}

\paragraph{Predictable $\sigma$-algebra.}

Fix a filtered probability space $(\Omega,\mathcal{F},(\mathcal{F}_t)_{t\in[0,T]},P)$, where
$\mathcal{F}_t$ denotes the $\sigma$-algebra of events observable up to time~$t$
(the \emph{filtration} of available information).

Define the \textbf{predictable $\sigma$-algebra} $\mathcal{P}_T$ on $\Omega \times [0,T]$ by
\[
    \mathcal{P}_T
    \;:=\;
    \sigma\!\Bigl(
        \bigl\{]s,t] \times F_s :\;
        0 \le s < t \le T,\;
        F_s \in \mathcal{F}_s
        \bigr\}
        \;\cup\;
        \bigl\{
        \{0\} \times F_0 :\;
        F_0 \in \mathcal{F}_0
        \bigr\}
    \Bigr).
\]

Equivalently,
\[
    \mathcal{P}_T
    \;=\;
    \sigma\!\bigl(
        \{(H_s)_{s \in [0,T]} :\;
        (H_s) \text{ is left-continuous and }
        (\mathcal{F}_t)\text{-adapted}
        \}
    \bigr).
\]
This is called the \textbf{predictable $\sigma$-algebra}.

\begin{remark}
The predictable $\sigma$-algebra $\mathcal{P}_T$ is generated by sets of the form $]s,t] \times F_s$ with $F_s \in \mathcal{F}_s$ (``stochastic rectangles''), together with $\{0\} \times F_0$.
Intuitively, a process is $\mathcal{P}_T$-measurable if and only if it is \textbf{left-continuous and adapted} (or, more precisely, a limit of such processes).
Predictability is a slightly stronger condition than being merely adapted: it requires the integrand at time~$t$ to be determined by information strictly \emph{before} time~$t$, rather than at time~$t$.
This is the correct measurability condition for integrands in It\^o's stochastic integral, ensuring that the integrand does not ``look into the future'' of the driving noise.
\end{remark}

\paragraph{The closure $\overline{\mathcal{E}}$.}

Then
\[
    \overline{\mathcal{E}}
    \;=\;
    L^2\!\bigl(\Omega_T,\;\mathcal{P}_T,\;P_T;\;L_2^0\bigr)
    \big/{\sim},
\]
where $\sim$ denotes identification up to classes of equivalence, and:
\begin{itemize}
    \item $\Omega_T := \Omega \times [0,T]$,
    \item $P_T := P \otimes dt$ \quad (product of the probability measure with Lebesgue measure on $[0,T]$),
    \item $U_0 := \sqrt{Q}\,U$,
    \item $L_2^0 := L_2(U_0,H)$
          is the space of Hilbert--Schmidt operators from $U_0$ to $H$.
\end{itemize}

\begin{remark}
The space $\overline{\mathcal{E}} = L^2(\Omega_T, \mathcal{P}_T, P_T; L_2^0)/{\sim}$ consists of all $\mathcal{P}_T$-measurable processes $\Phi \colon \Omega \times [0,T] \to L_2^0$ satisfying
\[
    \|\Phi\|_{\overline{\mathcal{E}}}^2
    \;:=\;
    E\!\left(\int_0^T
    \bigl\|\Phi(s) \circ \sqrt{Q}\,\bigr\|_{L_2(U,H)}^2\,ds
    \right)
    \;<\; \infty,
\]
with two processes identified if they agree $P_T$-almost everywhere.
The notation $L_2^0 = L_2(U_0, H)$ with $U_0=\sqrt{Q}\,U$ highlights that we only require $\Phi(s) \circ \sqrt{Q}$ (not $\Phi(s)$ itself) to be Hilbert--Schmidt (when $\Phi(s)$ is represented as an operator on~$U$).
This is natural because the Wiener--It\^o isometry involves $\|\Phi_s \circ \sqrt{Q}\|_{L_2(U,H)}$:
the noise enters through $\sqrt{Q}$, so the integrand need only be ``well-behaved'' on the range of~$\sqrt{Q}$.
\end{remark}

%% =============================
%% Step 4: Localization
%% =============================
\subsection*{Step 4: Localization}

As in the finite-dimensional setting, the stochastic integral $I$ can be defined on all
\[
    \Phi \colon \Omega_T \longrightarrow L_2^0,
    \qquad
    \text{adapted, } \mathcal{P}_T\text{-measurable},
\]
subject only to the condition
\[
    P\!\left(
        \int_0^T
        \bigl\|\Phi(s) \circ \sqrt{Q}\,\bigr\|_{L_2(U,H)}^2
        \,ds
        \;<\; \infty
    \right)
    \;=\; 1.
\]

\begin{remark}
In Steps~2--3, the stochastic integral was defined for integrands satisfying the \textbf{$L^2$-condition}:
$E\!\bigl(\int_0^T \|\Phi(s) \circ \sqrt{Q}\|_{L_2(U,H)}^2\,ds\bigr) < \infty$
(finiteness \emph{in expectation}).
The localization step weakens this to a \textbf{pathwise (almost sure)} condition:
$\int_0^T \|\Phi(s) \circ \sqrt{Q}\|_{L_2(U,H)}^2\,ds < \infty$ holds $P$-a.s.,
without requiring the expectation to be finite.
This extension is achieved by a \textbf{stopping-time argument}: one ``truncates'' $\Phi$ to make it $L^2$-integrable, defines the integral on each truncation, and verifies consistency.
The resulting integral is a \textbf{local martingale} (rather than a true martingale), exactly as in the scalar It\^o theory.
This is the largest natural domain for the stochastic integral.
\end{remark}

%% =============================
%% Properties of the stochastic integral
%% =============================
\subsection*{Properties of the stochastic integral}

%% =============================
%% (i) Linearity / composition with bounded operators
%% =============================
\paragraph{(i) Linearity (composition with bounded operators).}

For $L \in L(H, \widetilde{H})$ (a bounded linear operator between Hilbert spaces, with
$\widetilde{H}$ another Hilbert space), we have
\[
    L\!\left(\int_0^t \Phi(s)\,dW(s)\right)
    \;=\;
    \int_0^t L \circ \Phi(s)\,dW(s).
\]

\begin{remark}
This property says that the stochastic integral \textbf{commutes with bounded linear operators}.
Concretely: applying a bounded linear transformation $L \colon H \to \widetilde{H}$ to the integral $\int_0^t \Phi(s)\,dW(s) \in H$ is the same as integrating the composed process $L \circ \Phi(s) \colon U \to \widetilde{H}$.
The proof follows directly from the definition for elementary processes (where $L$ can be pulled through the finite sum) and extends to general integrands by the density/isometry argument.
This is the infinite-dimensional analogue of the scalar property $c \cdot \int_0^t \phi(s)\,dB_s = \int_0^t c\,\phi(s)\,dB_s$ for a constant $c \in \mathbb{R}$.
\end{remark}

%% =============================
%% (ii) Inner product with adapted process
%% =============================
\paragraph{(ii$'$) Inner product with a fixed vector.}

Let $h \in H$ be fixed. Then
\[
    \left\langle
        h,\;
        \int_0^t \Phi(s)\,dW(s)
    \right\rangle_{\!H}
    \;=\;
    \int_0^t \underbrace{h^*\!\circ\Phi(s)}_{=\;\widetilde{\Phi}_{h}(s)}\,dW_s,
\]
where
$\widetilde{\Phi}_{h}(s) := h^*\!\circ\Phi(s) \in L(U, \mathbb{R})$,
with $h^* = \langle h,\,\cdot\,\rangle_H \colon H \to \mathbb{R}$
denoting the linear functional induced by $h \in H$ via the Riesz representation.

\begin{remark}
Property (ii$'$) says that the fixed linear functional $\langle h,\,\cdot\,\rangle_H$ can be ``pulled inside'' the stochastic integral.
The operator $\widetilde{\Phi}_{h}(s) = \langle h,\,\Phi(s)(\cdot)\rangle_H : U \to \mathbb{R}$ maps $u \in U$ to the scalar $\langle h,\,\Phi(s)\,u\rangle_H$, so the right-hand side is a \textbf{real-valued} stochastic integral.
Equivalently, for each fixed $h$, the scalar process
$t \mapsto \langle h, \int_0^t \Phi(s)\,dW(s) \rangle_H$
is exactly the It\^o integral of the scalar integrand $u \mapsto \langle h, \Phi(s)u\rangle_H$:
\[
    \left\langle h, \int_0^t \Phi(s)\,dW(s) \right\rangle_H
    =
    \int_0^t h^*\Phi(s)\,dW(s).
\]
This follows first for elementary processes by linearity of the inner product and then for general $\Phi$ by the $L^2$-isometry/approximation argument.
The identification of $h$ with the functional $h^*=\langle h,\cdot\rangle_H$ is precisely the Riesz representation theorem in the Hilbert space~$H$.
For a time-varying (possibly random) adapted process $\hat{f}(t)$, the displayed identity is \emph{not} valid in general; one must use an It\^o product rule (integration by parts), which produces additional terms involving $d\hat{f}$ and quadratic covariation.
\end{remark}

%% =============================
%% (iii) BDG, (iv) quadratic variation, (v) regularity
%% =============================
\paragraph{(iii) Burkholder--Davis--Gundy (BDG) inequalities.}

For every $p \ge 1$ there exists a constant $C_p$ (depending only on $p$) such that for every integrand $\Phi \in \overline{\mathcal{E}}$,
\[
        E\Bigl(\sup_{0 \le t \le T} \big\|\int_0^t \Phi(s)\,dW(s)\big\|_H^p\Bigr)
        \le C_p\; E\Bigl(\Bigl(\int_0^T \|\Phi(s) \circ \sqrt{Q}\|_{L_2(U,H)}^2\,ds\Bigr)^{p/2}\Bigr).
\]

\begin{remark}
The BDG inequalities are the natural maximal estimates for martingales.  They follow from the scalar BDG applied to the canonical coordinate Brownian motions together with the Hilbert--Schmidt representation of the noise; alternatively one can reduce to the finite-dimensional case via projections and then pass to the limit using the isometry.  BDG is essential for pathwise regularity estimates and for controlling the supremum norm of stochastic integrals.
\end{remark}

\paragraph{(iv) Quadratic variation.}

Let $M_t := I(\Phi)_t = \int_0^t \Phi(s)\,dW(s)$.  Then the predictable quadratic variation of $M$ is given by
\[
        \langle M \rangle_t 
        \;=
        \;\int_0^t \|\Phi(s) \circ \sqrt{Q}\|_{L_2(U,H)}^2\,ds,
        \qquad t \in [0,T].
\]
Moreover the process $\|M_t\|_H^2 - \langle M \rangle_t$ is a real-valued martingale.

\begin{remark}
The identity for the quadratic variation is the infinite-dimensional analogue of
$\langle \int_0^\cdot \phi_s\,dB_s\rangle_t = \int_0^t \phi_s^2\,ds$.
It is proved first for step integrands by direct computation of cross terms,
then extended to general $\Phi$ by density and the Wiener--It\^o isometry.
The process $\|M_t\|_H^2 - \langle M\rangle_t$ is a martingale by the It\^o formula for
$\|M_t\|_H^2$ (or equivalently by applying the scalar isometry to each coordinate).
\end{remark}

\paragraph{(v) Regularity of the stochastic integral (Sobolev--Slobodetskii).}

Fix $\alpha \in (0,\tfrac12)$.  Define the fractional Sobolev (Slobodetskii) seminorm for an $H$-valued process $M$ by
\[
        [M]_{W^{\alpha,2}([0,T];H)}^2
        :=
        \int_0^T\int_0^T \frac{\|M_t - M_s\|_H^2}{|t-s|^{1+2\alpha}}\,ds\,dt.
\]
If $\Phi \in \overline{\mathcal{E}}$ satisfies the $L^2$-condition, then $M = I(\Phi)$ belongs to the space $L^2(\Omega;W^{\alpha,2}([0,T];H))$ and one has the estimate
\[
        E\bigl([M]_{W^{\alpha,2}}^2\bigr)
        \le C(\alpha)\; E\Bigl(\int_0^T \|\Phi(s) \circ \sqrt{Q}\|_{L_2(U,H)}^2\,ds\Bigr)
        \;<\; \infty.
\]

\begin{remark}
Recall that
\[
    W^{\alpha,2}([0,T];H)
    =
    \bigl\{f\in L^2(0,T;H): [f]_{W^{\alpha,2}([0,T];H)} < \infty\bigr\}.
\]
Let $M=I(\Phi)$.  By symmetry,
\[
    E\bigl([M]_{W^{\alpha,2}}^2\bigr)
    =
    2\int_0^T\int_t^T
    \frac{E\bigl(\|M_s-M_t\|_H^2\bigr)}{(s-t)^{1+2\alpha}}\,ds\,dt.
\]
For $0\le t<s\le T$ we have $M_s-M_t=\int_t^s \Phi(r)\,dW(r)$, hence by the It\^o isometry,
\[
    E\bigl(\|M_s-M_t\|_H^2\bigr)
    =
    E\Bigl(\int_t^s \|\Phi(r)\circ\sqrt{Q}\|_{L_2(U,H)}^2\,dr\Bigr).
\]
Insert this and apply Tonelli:
\begin{align*}
    E\bigl([M]_{W^{\alpha,2}}^2\bigr)
    &=
    2\,E\int_0^T\int_t^T\int_t^s
    \frac{\|\Phi(r)\circ\sqrt{Q}\|_{L_2(U,H)}^2}{(s-t)^{1+2\alpha}}\,dr\,ds\,dt \\
    &=
    2\,E\int_0^T \|\Phi(r)\circ\sqrt{Q}\|_{L_2(U,H)}^2
    \Bigl(\int_0^r\int_r^T (s-t)^{-1-2\alpha}\,ds\,dt\Bigr)\,dr.
\end{align*}
For each fixed $r\in[0,T]$, the kernel is finite when $\alpha<\tfrac12$ and satisfies
\[
    \int_0^r\int_r^T (s-t)^{-1-2\alpha}\,ds\,dt
    \le
    C(\alpha) < \infty,
\]
so we obtain
\[
    E\bigl([M]_{W^{\alpha,2}}^2\bigr)
    \le
    C(\alpha)\;E\Bigl(\int_0^T \|\Phi(r)\circ\sqrt{Q}\|_{L_2(U,H)}^2\,dr\Bigr).
\]
Together with $E\int_0^T \|M_t\|_H^2\,dt < \infty$ (from the It\^o isometry), this yields
$M\in L^2(\Omega;W^{\alpha,2}([0,T];H))$.
\end{remark}

%% =============================
%% Cylindrical Wiener process
%% =============================
\subsection*{Cylindrical Wiener process}

The $Q$-Wiener process
$W_t = \sum_{n\ge 1} \sqrt{\lambda_n}\,\beta_n(t)\,e_n$
converges in $L^2(\Omega;U)$ provided $\tr(Q) = \sum_n \lambda_n < \infty$.
In this case the stochastic integral
$\int_0^t \Phi(s)\,dW_s$ is well-defined for predictable $\Phi$ such that
\[
    E\int_0^T \|\Phi(s)\|_{L_2(U_0,H)}^2\,ds
    \;=\;
    E\int_0^T \|\Phi(s)\circ\sqrt{Q}\|_{L_2(U,H)}^2\,ds
    \;<\;\infty.
\]

When $Q$ is not trace class one can still make sense of the ``noise'' by working with the
\textbf{cylindrical Wiener process}
\[
    W_t
    \;=\;
    \sum_{n \ge 1} \beta_n(t)\,e_n,
\]
which does \emph{not} converge in $U$.  In that case the stochastic integral is defined via
\[
    \int_0^t \Phi(s)\,dW_s
    \;:=\;
    \sum_{n\ge 1}\int_0^t \Phi(s)e_n\,d\beta_n(s),
\]
and it is well-defined as an $H$-valued continuous local martingale whenever $\Phi$ is predictable and
\[
    \int_0^T \|\Phi(s)\|_{L_2(U,H)}^2\,ds < \infty \quad P\text{-a.s.}
\]
(in particular, if the above integral is finite in expectation).

\begin{remark}
The cylindrical Wiener process corresponds formally to $Q = \mathrm{Id}$
(all eigenvalues equal to~$1$).  In this case $\tr(Q) = \infty$, so the
series $\sum_n \beta_n(t)\,e_n$ does not converge in $U$.
However, the stochastic integral against it is still meaningful: since
$\sqrt{Q} = \mathrm{Id}$ and $U_0=U$, one needs $\Phi(s)$ itself to be
Hilbert--Schmidt from $U$ to~$H$, i.e.\ $\Phi(s) \in L_2(U,H)$.
Cylindrical Wiener processes arise naturally in practice (e.g.\ space--time
white noise on $L^2$), and the entire stochastic integration theory
developed above carries over to this setting without change, since the
integral was always defined in terms of $\Phi \circ \sqrt{Q}$.
\end{remark}

%% =============================
%% Stochastic evolution equations on Hilbert spaces
%% =============================
\section*{Stochastic evolution equations on Hilbert spaces}

Throughout this section we fix two separable (real) Hilbert spaces $U,H$ and a
$Q$-Wiener process $(W_t)_{t\ge 0}$ with respect to a filtration $(\mathcal{F}_t)_{t\ge 0}$
on~$U$ (or, when $Q$ is not trace class, a cylindrical Wiener process).

\paragraph{The (infinite-dimensional) SPDE / stochastic evolution equation.}
Consider the stochastic evolution equation
\begin{equation}\label{eq:spde}\tag{4.1}
		    dX_t
		    \;=\;
		    \bigl[A\,X_t + B(X_t)\bigr]\,dt
		    \;+\;
	    C(X_t)\,dW_t
	    \quad\text{in } H,
	    \qquad
	    X_0 = \xi,
\end{equation}
\begin{remark}
In \emph{(4.1)}, the unknown is an $H$-valued stochastic process
$X=(X_t)_{t\in[0,T]}$.  The term $A\,X_t$ is the \emph{linear drift}
(typically generated by a differential operator, e.g.\ a Laplacian-type
operator), while $B(X_t)$ is a (possibly nonlinear) \emph{drift/nonlinear
forcing} term.  The term $C(X_t)\,dW_t$ is the \emph{diffusion/noise} term:
for each state $x\in H$, $C(x)$ is an operator-valued coefficient (from the
noise space into $H$), and $W$ is the driving $Q$-Wiener (or cylindrical
Wiener) process.  Finally, $\xi$ is the initial condition at time $t=0$.
Thus \emph{(4.1)} is the infinite-dimensional analogue of a finite-dimensional
It\^o SDE. In this setting one often calls \emph{(4.1)} an SPDE (or stochastic
evolution equation), with $A$ encoding the spatial evolution operator and $C$
describing how the noise enters the system.
\end{remark}
subject to the following assumptions.

%% =============================
%% Assumptions (A1)--(A4)
%% =============================
\paragraph{Assumptions.}

\begin{itemize}
    \item[\textbf{(A1)}]
        $(A, D(A))$ generates a $C_0$-semigroup $(T_t)_{t \ge 0}$ on $H$.

    \item[\textbf{(A2)}]
        $B \colon H \to H$ is $\mathcal{B}(H)$-measurable.

    \item[\textbf{(A3)}]
        $C \colon H \to L_2(U_0, H)$ is \textbf{strongly continuous},
        i.e.\ for each fixed $u \in U_0$, the map
        \[
            x \;\longmapsto\; C(x)\,u,
            \qquad H \to H,
        \]
        is continuous.
        Here $U_0:=\sqrt{Q}\,U$ denotes the \emph{effective noise space} associated with the
        covariance operator $Q$ (the range of $\sqrt{Q}$, equipped with the natural Hilbert
        structure used in the stochastic integral).

    \item[\textbf{(A4)}]
        $\xi$ is $H$-valued and $\mathcal{F}_0$-measurable.
\end{itemize}

\begin{remark}
\textbf{(A1)} says that the linear part $A$ of the drift generates a
strongly continuous (i.e.\ $C_0$-) semigroup.  This is the standard
framework for abstract parabolic and hyperbolic evolution equations
(e.g.\ the Laplacian on $L^2$ generates the heat semigroup).

\textbf{(A2)} is a minimal measurability assumption on the nonlinear
drift~$B$.  In applications one typically imposes Lipschitz or
local Lipschitz conditions on~$B$ to obtain existence and uniqueness;
here we state the weakest structural requirement.
In particular, if $X$ is an $H$-valued adapted/measurable process, then
$B(X_t)$ is again $H$-valued measurable for each $t$, so the drift term
can be interpreted as a Bochner integral (under suitable integrability).

\textbf{(A3)} requires the diffusion coefficient $C$ to be
\emph{strongly} continuous (continuity of $C(x)\,u$ in $x$ for each
fixed $u$), rather than norm-continuous.  Strong continuity is
strictly weaker than norm (operator-norm) continuity and suffices
for the measurability properties needed in the stochastic integral.

\textbf{(A4)} is the usual initial-condition assumption:
$\xi$ is an $H$-valued random variable known at time~$0$.
\end{remark}

%% =============================
%% C_0-semigroup
%% =============================
\subsection*{$C_0$-semigroups}

A family of bounded linear operators $(T_t)_{t \ge 0}$ on $H$
is called a \textbf{$C_0$-semigroup} if:
\begin{enumerate}
    \item $T_0 = \mathrm{Id}$,
    \item $T_t \circ T_s = T_{t+s}$ for all $t, s \ge 0$ \quad (semigroup property),
    \item $t \mapsto T_t\,u$ is continuous for every $u \in H$
          \quad (strong continuity).
\end{enumerate}

The \textbf{generator} of the semigroup is the (generally unbounded)
operator $A$ defined by
\[
    A\,h
    \;:=\;
    \lim_{t \to 0^+}
    \frac{T_t\,h - T_0\,h}{t}
    \;=\;
    \lim_{t \to 0^+}
    \frac{T_t\,h - h}{t},
\]
with domain $D(A) := \{h \in H : \text{the limit above exists in } H\}$.

\begin{remark}
The generator $A$ encodes the ``infinitesimal'' behaviour of the
semigroup: $T_t \approx \mathrm{Id} + t\,A$ for small~$t$.
This explains the terminology \emph{infinitesimal generator}: for $h\in D(A)$,
the limit $Ah=\lim_{t\to 0^+}(T_th-h)/t$ is the time derivative of
$t\mapsto T_t h$ at $t=0$, so $A$ describes the first-order effect of an
infinitesimal time step.
For the heat equation on $H = L^2(\mathbb{R}^d)$,
$A = \tfrac{1}{2}\Delta$ and $T_t$ is the heat semigroup
(convolution with the Gaussian kernel).  More precisely,
for $t>0$ let
\[
    G_t(x) := (2\pi t)^{-d/2}\exp\!\bigl(-|x|^2/(2t)\bigr),
    \qquad x\in\mathbb{R}^d,
\]
then for $u_0\in L^2(\mathbb{R}^d)$ one has
\[
    (T_t u_0)(x) = (G_t*u_0)(x) := \int_{\mathbb{R}^d} G_t(x-y)\,u_0(y)\,dy.
\]
The semigroup property $T_tT_s=T_{t+s}$ corresponds to the Gaussian convolution identity
$G_t*G_s = G_{t+s}$.
In general $A$ is an
\emph{unbounded} operator (defined only on a dense domain $D(A)
\subset H$), and the Hille--Yosida theorem characterises exactly
which operators generate $C_0$-semigroups.

The semigroup property $T_t \circ T_s = T_{t+s}$ is the operator
analogue of $e^{at}\,e^{as} = e^{a(t+s)}$; indeed when $A$ is bounded
one simply has $T_t = e^{tA}$.
\end{remark}

%% =============================
%% Notions of solution
%% =============================
\subsection*{Notions of solution}

We discuss three standard notions of solutions for \eqref{eq:spde}.

\paragraph{Mild solution.}
An $H$-valued predictable process $(X_t)_{t\in[0,T]}$ is called a
\textbf{mild solution} to \eqref{eq:spde} if for all $t\in[0,T]$,
\begin{equation}\tag{4.2}
    X_t
    \;=\;
    T_t\,\xi
    \;+
    \int_0^t T_{t-s}\,B\bigl(X_s\bigr)\,ds
    \;+
    \int_0^t T_{t-s}\,C\bigl(X_s\bigr)\,dW_s,
    \qquad P\text{-a.s.}
\end{equation}

\begin{remark}
The mild formulation replaces the (potentially ill-defined) term $A\,X_t$ by
the semigroup $T_{t-s}$ (variation of constants / Duhamel principle).
Intuitively, $T_{t-s}$ is the linear flow generated by $A$ from time $s$ to time $t$:
it propagates a contribution created at time $s$ forward to time $t$ under the linear dynamics.
Thus the forcing terms $B(X_s)\,ds$ and $C(X_s)\,dW_s$ injected at time $s$ appear at time $t$
as $T_{t-s}B(X_s)\,ds$ and $T_{t-s}C(X_s)\,dW_s$.
This is crucial because solutions typically do not lie in $D(A)$,
so the expression $A\,X_t$ may not make sense pointwise.
The two integrals in \emph{(4.2)} must be well-defined: the Bochner integral
for $T_{t-s}B(X_s)$ and the stochastic integral for $T_{t-s}C(X_s)$.
\end{remark}

\paragraph{Strong (analytical) solution.}
An $H$-valued predictable process $(X_t)$ is called a
\textbf{(analytically) strong solution} if
\begin{itemize}
    \item $X_t \in D(A)$ for a.e.\ $t$ (and $P$-a.s.),
    \item $\int_0^T \|A\,X_s\|_H\,ds < \infty$ and the integrals below exist,
\end{itemize}
and for all $t\in[0,T]$,
\[
    X_t
    \;=\;
    \xi
    \;+
    \int_0^t \bigl[A\,X_s + B(X_s)\bigr] \,ds
    \;+
    \int_0^t C(X_s)\,dW_s,
    \qquad P\text{-a.s.}
\]

\begin{remark}
The term \emph{strong solution} is used in two different senses in the
stochastic evolution equation literature: (i)\ strong in the
probabilistic sense (solution adapted to the given Wiener process), and
(ii)\ strong in the analytical sense (solution
possesses enough spatial regularity so that $A\,X_t$ is meaningful).
Here we mean the \textbf{analytical} notion.
\end{remark}

\paragraph{Weak (analytical) solution.}
An $H$-valued predictable process $(X_t)$ is called an
\textbf{(analytically) weak solution} if for every test vector
$\varphi \in D(A^*)$ and all $t\in[0,T]$,
\[
    \langle X_t,\varphi\rangle_H
    \;=\;
    \langle \xi,\varphi\rangle_H
    \;+
    \int_0^t \langle X_s, A^*\varphi\rangle_H\,ds
    \;+
    \int_0^t \langle B(X_s),\varphi\rangle_H\,ds
    \;+
    \int_0^t \langle \varphi, C(X_s)\,dW_s\rangle_H,
    \qquad P\text{-a.s.}
\]
where $(A^*,D(A^*))$ denotes the adjoint operator of $A$.

\begin{remark}
In the weak formulation, one tests the equation against $\varphi \in D(A^*)$
so that $\langle A\,X_s,\varphi\rangle_H$ can be interpreted as
$\langle X_s, A^*\varphi\rangle_H$ without requiring $X_s\in D(A)$.
This parallels the classical PDE notion of weak solution.
\end{remark}

%% =============================
%% Existence and uniqueness of mild solutions
%% =============================
\subsection*{Existence and uniqueness of mild solutions}

We impose the following conditions on $B$ and $C$.

\paragraph{(H.1) Lipschitz condition.}
There exists $L>0$ such that for all $x,y\in H$,
\[
    \|B(x)-B(y)\|_H^2
    \;+
    \|\bigl(C(x)-C(y)\bigr)\circ\sqrt{Q}\|_{L_2(U,H)}^2
    \;\le\;
    L\,\|x-y\|_H^2.
\]

\paragraph{(H.2) Linear growth.}
There exists $M>0$ such that for all $x\in H$,
\[
    \|B(x)\|_H^2
    \;+
    \|C(x)\circ\sqrt{Q}\|_{L_2(U,H)}^2
    \;\le\;
    M\,\bigl(1+\|x\|_H^2\bigr).
\]

\begin{theorem}[Existence and uniqueness of a mild solution]
Assume \textbf{(A1)}--\textbf{(A4)} and \textbf{(H.1)}--\textbf{(H.2)}.  In addition, assume that
$E\|\xi\|_H^p<\infty$ for some $p\ge 2$ (equivalently $\xi\in L^p(\Omega;H)$).
Then there exists a unique mild solution $(X_t)_{t\in[0,T]}$ of \eqref{eq:spde}.
Moreover,
\[
    P\!\left(\int_0^T \|X_s\|_H^2\,ds < \infty\right) = 1.
\]
\end{theorem}

\begin{remark}
The proof is typically based on a fixed-point argument (Picard iteration)
on the space of predictable processes, using the Wiener--It\^o isometry and
BDG inequalities to control the stochastic convolution
$\int_0^t T_{t-s}C(X_s)\,dW_s$.
The Lipschitz assumption \textbf{(H.1)} yields contraction (locally in time,
then global by iteration), while \textbf{(H.2)} ensures integrability.
The assumption $E\|\xi\|_H^p<\infty$ is an \emph{integrability hypothesis} on the initial
condition, needed because we work in the $L^p$-based space $H_p$ and must control
the term $T_t\xi$ in $L^p$.  It is satisfied, for instance, if $\xi$ is deterministic, bounded,
or more generally $\xi\in L^p(\Omega;H)$.
\end{remark}

\subsection*{Proof via Banach's fixed point theorem (details)}

Fix $p\ge 2$ and define the Banach space
\[
    H_p
    :=
    \Bigl\{Y:[0,T]\times\Omega\to H:\ Y\text{ is predictable and }
    \|Y\|_{H_p}<\infty\Bigr\},
    \qquad
    \|Y\|_{H_p}
    :=
    \sup_{t\in[0,T]}\bigl(E\|Y_t\|_H^p\bigr)^{1/p}.
\]

For $Y\in H_p$, define the map $K$ by
\begin{equation}\label{eq:K-map}
    (KY)(t)
    :=
    T_t\,\xi
    +
    \int_0^t T_{t-s}\,B\bigl(Y_s\bigr)\,ds
    +
    \int_0^t T_{t-s}\,C\bigl(Y_s\bigr)\,dW_s,
    \qquad t\in[0,T].
\end{equation}
Write
\[
    (KY)(t) = T_t\xi + K_1(Y)(t) + K_2(Y)(t)
\]
with
\[
    K_1(Y)(t):=\int_0^t T_{t-s}B(Y_s)\,ds,
    \qquad
    K_2(Y)(t):=\int_0^t T_{t-s}C(Y_s)\,dW_s.
\]
Let $M_T:=\sup_{r\in[0,T]}\|T_r\|_{L(H)}<\infty$.

\paragraph{Step 1: $K$ maps $H_p$ into itself.}
Let $Y\in H_p$.  Since $Y$ is predictable and $B,C$ are measurable,
the processes $B(Y)$ and $C(Y)$ are predictable, and therefore
$KY$ is predictable as well.  It remains to show that
$\sup_{t\in[0,T]}E\|KY(t)\|_H^p<\infty$.

First, we use the elementary inequality
\begin{equation}\label{eq:sum-p-ineq}
    \|a+b+c\|_H^p
    \le
    3^{p-1}\bigl(\|a\|_H^p+\|b\|_H^p+\|c\|_H^p\bigr),
    \qquad a,b,c\in H,
\end{equation}
which follows from Jensen's inequality applied to the convex function $u\mapsto \|u\|_H^p$.
Thus, for each $t\in[0,T]$,
\[
    E\|KY(t)\|_H^p
    \le
    3^{p-1}\Bigl(E\|T_t\xi\|_H^p + E\|K_1(Y)(t)\|_H^p + E\|K_2(Y)(t)\|_H^p\Bigr).
\]

\paragraph{Drift term.}
By $\|T_{t-s}\|_{L(H)}\le M_T$,
\[
    \|K_1(Y)(t)\|_H
    =
    \Bigl\|\int_0^t T_{t-s}B(Y_s)\,ds\Bigr\|_H
    \le
    \int_0^t \|T_{t-s}\|_{L(H)}\,\|B(Y_s)\|_H\,ds
    \le
    M_T\int_0^t \|B(Y_s)\|_H\,ds.
\]
For $p\ge 1$, H\"older's inequality gives
\[
    \Bigl(\int_0^t f(s)\,ds\Bigr)^p
    \le
    t^{p-1}\int_0^t f(s)^p\,ds,
    \qquad f\ge 0,
\]
and therefore,
\begin{equation}\label{eq:k1-bound}
    E\|K_1(Y)(t)\|_H^p
    \le
    M_T^p\,t^{p-1}\int_0^t E\|B(Y_s)\|_H^p\,ds.
\end{equation}
Since the integrand is nonnegative, Tonelli yields
\[
    E\int_0^t \|B(Y_s)\|_H^p\,ds
    =
    \int_0^t E\|B(Y_s)\|_H^p\,ds.
\]
Moreover, by \textbf{(H.2)} and $(1+\|x\|_H^2)^{p/2}\le c_p(1+\|x\|_H^p)$,
\[
    E\|B(Y_s)\|_H^p
    \le
    C(p,M)\bigl(1+E\|Y_s\|_H^p\bigr)
    \le
    C(p,M)\bigl(1+\|Y\|_{H_p}^p\bigr).
\]
In particular, combining with \eqref{eq:k1-bound} we obtain the simplified estimate
\[
    E\|K_1(Y)(t)\|_H^p
    \le
    C(p,M,M_T)\,t^{p}\,\bigl(1+\|Y\|_{H_p}^p\bigr)
    \le
    C(p,M,M_T)\,T^{p}\,\bigl(1+\|Y\|_{H_p}^p\bigr).
\]

\paragraph{Diffusion term.}
Fix $t\in[0,T]$ and consider the $H$-valued martingale
\[
    N_r := \int_0^r T_{t-s}C(Y_s)\,dW_s,
    \qquad r\in[0,t].
\]
Then $K_2(Y)(t)=N_t$, and BDG (in fixed time form, obtained from the maximal inequality) gives
\[
    E\|K_2(Y)(t)\|_H^p
    \le
    C_p\,E\Bigl(\int_0^t \|T_{t-s}C(Y_s)\circ\sqrt{Q}\|_{L_2(U,H)}^2\,ds\Bigr)^{p/2}.
\]
Using $\|T_{t-s}\|_{L(H)}\le M_T$ and $\|T_{t-s}S\|_{L_2(U,H)}\le \|T_{t-s}\|_{L(H)}\,\|S\|_{L_2(U,H)}$,
we obtain
\[
    E\|K_2(Y)(t)\|_H^p
    \le
    C_p\,M_T^p\,E\Bigl(\int_0^t \|C(Y_s)\circ\sqrt{Q}\|_{L_2(U,H)}^2\,ds\Bigr)^{p/2}.
\]
Since $p/2\ge 1$, the same H\"older estimate (with exponent $p/2$) yields
\[
    \Bigl(\int_0^t g(s)\,ds\Bigr)^{p/2}
    \le
    t^{\frac p2-1}\int_0^t g(s)^{p/2}\,ds,
    \qquad g\ge 0,
\]
and with $g(s)=\|C(Y_s)\circ\sqrt{Q}\|_{L_2(U,H)}^2$ we get
\begin{equation}\label{eq:k2-bound}
    E\|K_2(Y)(t)\|_H^p
    \le
    C_p\,M_T^p\,t^{\frac p2-1}\int_0^t E\|C(Y_s)\circ\sqrt{Q}\|_{L_2(U,H)}^p\,ds.
\end{equation}

Since the integrand is nonnegative, Tonelli gives
\[
    E\int_0^t \|C(Y_s)\circ\sqrt{Q}\|_{L_2(U,H)}^p\,ds
    =
    \int_0^t E\|C(Y_s)\circ\sqrt{Q}\|_{L_2(U,H)}^p\,ds.
\]
Moreover, by \textbf{(H.2)} and $(1+\|x\|_H^2)^{p/2}\le c_p(1+\|x\|_H^p)$,
\[
    E\|C(Y_s)\circ\sqrt{Q}\|_{L_2(U,H)}^p
    \le
    C(p,M)\bigl(1+E\|Y_s\|_H^p\bigr)
    \le
    C(p,M)\bigl(1+\|Y\|_{H_p}^p\bigr).
\]
In particular, combining with \eqref{eq:k2-bound} we obtain the simplified estimate
\[
    E\|K_2(Y)(t)\|_H^p
    \le
    C(p,M,M_T)\,t^{p/2}\,\bigl(1+\|Y\|_{H_p}^p\bigr)
    \le
    C(p,M,M_T)\,T^{p/2}\,\bigl(1+\|Y\|_{H_p}^p\bigr).
\]
Equivalently, for fixed $T>0$ one may write
\[
    E\|K_2(Y)(t)\|_H^p
    \le
    M(T)\,T^{(p+1)/2}\,\bigl(1+\|Y\|_{H_p}^p\bigr),
\]
for a constant $M(T)$ depending on $(p,M,M_T,T)$ (absorbing $T^{-1/2}$ into the constant).

Combining the bounds for $T_t\xi$, $K_1(Y)(t)$, and $K_2(Y)(t)$ and taking $\sup_{t\in[0,T]}$, we obtain
\[
    \sup_{t\in[0,T]}E\|KY(t)\|_H^p
    \le
    C\Bigl(E\|\xi\|_H^p + 1 + \|Y\|_{H_p}^p\Bigr),
\]
for a constant $C=C(p,T,M,M_T)$ (absorbing the factors $T^p$ and $T^{p/2}$).
Hence $KY\in H_p$ whenever $Y\in H_p$.

\paragraph{Step 2: Contraction estimate (small time).}
Let $Y_1,Y_2\in H_p$ and $t\in[0,T]$.  The semigroup term cancels, so
\[
    (KY_1)(t)-(KY_2)(t)
    =
    \bigl(K_1(Y_1)(t)-K_1(Y_2)(t)\bigr)
    +
    \bigl(K_2(Y_1)(t)-K_2(Y_2)(t)\bigr).
\]
We will use Jensen's inequality (applied to the convex function $u\mapsto \|u\|_H^p$):
\begin{equation}\label{eq:two-sum-p-ineq}
    \|u+v\|_H^p \le 2^{p-1}\bigl(\|u\|_H^p+\|v\|_H^p\bigr),
    \qquad u,v\in H.
\end{equation}

\paragraph{Drift difference.}
By definition,
\[
    K_1(Y_1)(t)-K_1(Y_2)(t)
    =
    \int_0^t T_{t-s}\bigl(B(Y_{1,s})-B(Y_{2,s})\bigr)\,ds.
\]
Thus, using $\|T_{t-s}\|\le M_T$, H\"older in time and Tonelli as in Step~1,
\[
\begin{aligned}
    E\|K_1(Y_1)(t)-K_1(Y_2)(t)\|_H^p
    &\le
    M_T^p\,t^{p-1}\int_0^t E\|B(Y_{1,s})-B(Y_{2,s})\|_H^p\,ds.
\end{aligned}
\]
From \textbf{(H.1)} we have $\|B(x)-B(y)\|_H^2\le L\|x-y\|_H^2$; raising to $p/2$ yields
\[
    \|B(x)-B(y)\|_H^p\le L^{p/2}\|x-y\|_H^p.
\]
Hence, for each $s$,
\[
    E\|B(Y_{1,s})-B(Y_{2,s})\|_H^p
    \le
    L^{p/2}\,E\|Y_{1,s}-Y_{2,s}\|_H^p
    \le
    L^{p/2}\,\|Y_1-Y_2\|_{H_p}^p,
\]
and therefore
\begin{equation}\label{eq:k1-diff-bound}
    E\|K_1(Y_1)(t)-K_1(Y_2)(t)\|_H^p
    \le
    M_T^p\,L^{p/2}\,t^{p}\,\|Y_1-Y_2\|_{H_p}^p.
\end{equation}

\paragraph{Diffusion difference.}
Similarly,
\[
    K_2(Y_1)(t)-K_2(Y_2)(t)
    =
    \int_0^t T_{t-s}\bigl(C(Y_{1,s})-C(Y_{2,s})\bigr)\,dW_s.
\]
By BDG (fixed time form), $\|T_{t-s}\|\le M_T$ and the same H\"older estimate in time,
\[
\begin{aligned}
    E\|K_2(Y_1)(t)-K_2(Y_2)(t)\|_H^p
    &\le
    C_p\,M_T^p\,t^{\frac p2-1}\int_0^t
    E\|\bigl(C(Y_{1,s})-C(Y_{2,s})\bigr)\circ\sqrt{Q}\|_{L_2(U,H)}^p\,ds.
\end{aligned}
\]
Using \textbf{(H.1)} in the form
$\|((C(x)-C(y))\circ\sqrt{Q})\|_{L_2(U,H)}^p\le L^{p/2}\|x-y\|_H^p$, we obtain
\begin{equation}\label{eq:k2-diff-bound}
    E\|K_2(Y_1)(t)-K_2(Y_2)(t)\|_H^p
    \le
    C_p\,M_T^p\,L^{p/2}\,t^{p/2}\,\|Y_1-Y_2\|_{H_p}^p.
\end{equation}

\paragraph{Combine and take $\sup_{t\in[0,T]}$.}
By \eqref{eq:two-sum-p-ineq}, \eqref{eq:k1-diff-bound}, and \eqref{eq:k2-diff-bound},
for each $t\in[0,T]$,
\[
    E\|(KY_1)(t)-(KY_2)(t)\|_H^p
    \le
    \widetilde{C}\,\bigl(t^{p}+t^{p/2}\bigr)\,\|Y_1-Y_2\|_{H_p}^p,
\]
for $\widetilde{C}=\widetilde{C}(p,L,M_T)$.
Taking $\sup_{t\in[0,T]}$ yields
\[
    \|KY_1-KY_2\|_{H_p}^p
    \le
    \widetilde{C}\,\bigl(T^{p}+T^{p/2}\bigr)\,\|Y_1-Y_2\|_{H_p}^p.
\]
Choosing $T_0>0$ such that $\widetilde{C}(T_0^{p}+T_0^{p/2})<1$, the map $K$
is a strict contraction on $H_p$ over $[0,T_0]$.

\paragraph{Step 3: Local fixed point and global solution.}
Since $H_p$ is complete and $K$ is a strict contraction on $[0,T_0]$,
Banach's fixed point theorem yields a unique fixed point $X\in H_p$ such that
$X=KX$ on $[0,T_0]$.  By definition of $K$, this is exactly the mild
formulation \emph{(4.2)} on $[0,T_0]$.
In particular, for the given initial condition $\xi$, there is a unique mild solution on the
time interval $[0,T_0]$.

To extend to $[0,T]$, one concatenates the local solutions.  Set $t_0:=0$ and
$t_{n+1}:=\min\{t_n+T_0,T\}$.  Assuming $X$ is constructed on $[0,t_n]$, define the next
segment on $[t_n,t_{n+1}]$ as the fixed point of the corresponding map with initial data
$X_{t_n}$, namely
\[
    X_t
    =
    T_{t-t_n}X_{t_n}
    +
    \int_{t_n}^t T_{t-s}B(X_s)\,ds
    +
    \int_{t_n}^t T_{t-s}C(X_s)\,dW_s,
    \qquad t\in[t_n,t_{n+1}].
\]
The contraction estimate depends only on the interval length (at most $T_0$) and the constants
in \textbf{(H.1)}--\textbf{(H.2)} and $M_T$, so the same argument applies on each interval
$[t_n,t_{n+1}]$.  Hence there is a unique mild solution on each such interval given the
left-endpoint value $X_{t_n}$, and by induction this yields existence and uniqueness on the
full horizon $[0,T]$.

\paragraph{Step 4: Moment bound and continuity.}
We briefly indicate the ingredients for the a priori moment bound.  Starting from the mild
formulation $X=KX$ and taking $\sup_{t\in[0,T]}$, one estimates:
(i)\ $\sup_t\|T_t\xi\|_H\le M_T\|\xi\|_H$ by boundedness of the semigroup,
(ii)\ the drift term by H\"older in time together with \textbf{(H.2)}, and
(iii)\ the stochastic term by a maximal BDG/Doob estimate together with $\|T_{t-s}\|\le M_T$
and \textbf{(H.2)}.  This yields a bound of the form
\[
    E\Bigl(\sup_{t\in[0,T]}\|X_t\|_H^p\Bigr)
    \le
    C_{p,T}\,\bigl(1+E\|\xi\|_H^p\bigr),
\]
with $C_{p,T}$ depending on $(p,T,L,M,M_T)$.
In particular, since $p\ge 2$,
\begin{align*}
    E\int_0^T \|X_s\|_H^2\,ds
    &=
    \int_0^T E\|X_s\|_H^2\,ds
    \\
    &\le
    \int_0^T \bigl(E\|X_s\|_H^p\bigr)^{2/p}\,ds
    \\
    &\le
    T\Bigl(\sup_{t\in[0,T]}E\|X_t\|_H^p\Bigr)^{2/p}
    \\
    &\le
    T\Bigl(E\sup_{t\in[0,T]}\|X_t\|_H^p\Bigr)^{2/p}
    <\infty,
\end{align*}
where we used Tonelli, the inequality $\|Z\|_{L^2}\le \|Z\|_{L^p}$ for $p\ge 2$, and the moment bound above.
In particular,
\[
    T\Bigl(E\sup_{t\in[0,T]}\|X_t\|_H^p\Bigr)^{2/p}
    \le
    T\,\Bigl(C_{p,T}\bigl(1+E\|\xi\|_H^p\bigr)\Bigr)^{2/p}
    <\infty.
\]
Since the random variable $\int_0^T \|X_s\|_H^2\,ds$ is nonnegative, finiteness of its expectation implies
\[
    P\!\left(\int_0^T \|X_s\|_H^2\,ds < \infty\right)=1,
\]
which is the integrability statement in the theorem (indeed, if it were infinite with positive probability,
then its expectation would be infinite).
Finally, $X$ admits a continuous modification (see the remark below).

\begin{remark}
\textbf{Why we care about a continuous modification.}
The process $X$ is initially constructed as a family of random variables $(X_t)_{t\in[0,T]}$,
each defined only up to changes on a null set that may depend on $t$.
A \emph{continuous modification} is another version $\widetilde X$ such that
$P(\widetilde X_t=X_t\ \text{for all }t)=1$ and, in addition, the sample paths
$t\mapsto \widetilde X_t(\omega)$ are continuous.
We care because many arguments in SPDE theory (and similarly in SDE theory) are \emph{pathwise} and require an actual trajectory
with good time-regularity: it lets us talk meaningfully about quantities such as
$\sup_{t\in[0,T]}\|X_t(\omega)\|_H$, sample-path H\"older regularity, and hitting/exit times
(which are stopping times when paths are continuous).
It also matches the usual definition of a solution as an $H$-valued process with continuous paths.

In the mild formula \eqref{eq:K-map}, each term has a continuous modification:
$t\mapsto T_t\xi$ is continuous by strong continuity of $(T_t)$,
the drift is a time-integral (hence continuous in $t$),
and the stochastic convolution admits a continuous modification (e.g.\ via BDG for increments together
with the regularity estimates from the integration theory above).
\end{remark}

\begin{remark}
\textbf{Intuition for Steps~1--4.}
\begin{itemize}
    \item Step~1 (define one Picard step): think of $K$ as the map ``plug a guess $Y$ into
    the right-hand side of the mild equation and compute the next iterate.''  To make this meaningful
    we must check two things: (i)\ measurability/predictability, so the stochastic integral is defined,
    and (ii)\ integrability, so the drift and diffusion terms do not blow up in $L^p$.  The growth bound
    \textbf{(H.2)} is exactly what turns $\|B(Y_s)\|$ and $\|C(Y_s)\circ\sqrt Q\|$ into something controlled
    by $1+\|Y_s\|$, allowing the estimates in $H_p$.
    \item Step~2 (make the Picard step contract): compare two guesses $Y_1,Y_2$ and estimate the size of
    $(KY_1)(t)-(KY_2)(t)$.  The Lipschitz bound \textbf{(H.1)} says that changing the
    input by $\|Y_1-Y_2\|$ changes the coefficients by at most a constant times $\|Y_1-Y_2\|$.  What makes
    this a \emph{contraction} is small time: the drift term involves an integral in time (so it brings in a
    factor of order $t$), while the stochastic term behaves like a Brownian integral (typical size of order
    $\sqrt{t}$).  After taking $p$-moments, these become factors $t^{p}$ and $t^{p/2}$, so for $T$ small
    enough the overall multiplicative factor is $<1$.
    \item Step~3 (cover the full time interval): Banach's theorem gives one unique fixed point on
    $[0,T_0]$, i.e.\ one unique solution path segment starting from $\xi$.  To go beyond $T_0$ you do not
    create independent solutions; you \emph{restart} with initial data equal to the endpoint value already
    produced, $X_{T_0}$, and apply the same local argument on $[T_0,2T_0]$, and so on.  This produces one
    single global path on $[0,T]$.
    \item Step~4 (turn estimates into path properties): once you have $L^p$ bounds on $\sup_{t\le T}\|X_t\|$,
    you can deduce $E\int_0^T\|X_s\|_H^2ds<\infty$, hence $\int_0^T\|X_s\|_H^2ds<\infty$ holds $P$-a.s.  The
    mild formula also shows that $X$ has a continuous modification: the semigroup term is continuous, the drift
    is a time-integral, and the stochastic convolution admits a continuous modification.
\end{itemize}
\end{remark}

\end{document}
